\documentclass[11pt,a4paper]{article}
% preamble.tex — estilo común de todos los apuntes Froth.
% Cada apunte hace % preamble.tex — estilo común de todos los apuntes Froth.
% Cada apunte hace % preamble.tex — estilo común de todos los apuntes Froth.
% Cada apunte hace \input{preamble} justo después de \documentclass.
\usepackage[margin=2.4cm]{geometry}
\usepackage{xcolor}
\definecolor{litblue}{RGB}{31,79,121}
\definecolor{litgreen}{RGB}{27,94,32}
\definecolor{litorange}{RGB}{230,81,0}
\usepackage{parskip}
\usepackage{titlesec}
\titleformat{\section}{\Large\bfseries\color{litblue}}{\thesection}{0.6em}{}
\titleformat{\subsection}{\large\bfseries\color{litblue}}{\thesubsection}{0.6em}{}
\usepackage{tcolorbox}
\usepackage{fancyhdr}
\pagestyle{fancy}
\fancyhf{}
\fancyhead[L]{\small\color{litblue}Froth — Apuntes de ML}
\fancyhead[R]{\small\thepage}
\renewcommand{\headrulewidth}{0.4pt}
\usepackage[colorlinks=true,linkcolor=litblue,urlcolor=litblue]{hyperref}

% Cabecera de cada apunte: \cabecera{numero}{titulo}
\newcommand{\cabecera}[2]{%
  \begin{tcolorbox}[colback=litblue,coltext=white,colframe=litblue,arc=2mm]
    {\Large\bfseries Apunte #1 — #2}\\[3pt]
    {\small Proyecto Froth · aprender Machine Learning construyendo}
  \end{tcolorbox}\vspace{0.8em}}

% Cajas temáticas
\newtcolorbox{concepto}[1]{colback=blue!4,colframe=litblue,title={Concepto: #1},arc=1.5mm}
\newtcolorbox{practica}{colback=green!5,colframe=litgreen,title={Pruébalo tú (teclado en mano)},arc=1.5mm}
\newtcolorbox{ojo}{colback=orange!8,colframe=litorange,title={Ojo / error típico},arc=1.5mm}
\newtcolorbox{resumen}{colback=gray!8,colframe=gray!60,title={En una frase},arc=1.5mm}

\newcommand{\codigo}[1]{\texttt{#1}}
 justo después de \documentclass.
\usepackage[margin=2.4cm]{geometry}
\usepackage{xcolor}
\definecolor{litblue}{RGB}{31,79,121}
\definecolor{litgreen}{RGB}{27,94,32}
\definecolor{litorange}{RGB}{230,81,0}
\usepackage{parskip}
\usepackage{titlesec}
\titleformat{\section}{\Large\bfseries\color{litblue}}{\thesection}{0.6em}{}
\titleformat{\subsection}{\large\bfseries\color{litblue}}{\thesubsection}{0.6em}{}
\usepackage{tcolorbox}
\usepackage{fancyhdr}
\pagestyle{fancy}
\fancyhf{}
\fancyhead[L]{\small\color{litblue}Froth — Apuntes de ML}
\fancyhead[R]{\small\thepage}
\renewcommand{\headrulewidth}{0.4pt}
\usepackage[colorlinks=true,linkcolor=litblue,urlcolor=litblue]{hyperref}

% Cabecera de cada apunte: \cabecera{numero}{titulo}
\newcommand{\cabecera}[2]{%
  \begin{tcolorbox}[colback=litblue,coltext=white,colframe=litblue,arc=2mm]
    {\Large\bfseries Apunte #1 — #2}\\[3pt]
    {\small Proyecto Froth · aprender Machine Learning construyendo}
  \end{tcolorbox}\vspace{0.8em}}

% Cajas temáticas
\newtcolorbox{concepto}[1]{colback=blue!4,colframe=litblue,title={Concepto: #1},arc=1.5mm}
\newtcolorbox{practica}{colback=green!5,colframe=litgreen,title={Pruébalo tú (teclado en mano)},arc=1.5mm}
\newtcolorbox{ojo}{colback=orange!8,colframe=litorange,title={Ojo / error típico},arc=1.5mm}
\newtcolorbox{resumen}{colback=gray!8,colframe=gray!60,title={En una frase},arc=1.5mm}

\newcommand{\codigo}[1]{\texttt{#1}}
 justo después de \documentclass.
\usepackage[margin=2.4cm]{geometry}
\usepackage{xcolor}
\definecolor{litblue}{RGB}{31,79,121}
\definecolor{litgreen}{RGB}{27,94,32}
\definecolor{litorange}{RGB}{230,81,0}
\usepackage{parskip}
\usepackage{titlesec}
\titleformat{\section}{\Large\bfseries\color{litblue}}{\thesection}{0.6em}{}
\titleformat{\subsection}{\large\bfseries\color{litblue}}{\thesubsection}{0.6em}{}
\usepackage{tcolorbox}
\usepackage{fancyhdr}
\pagestyle{fancy}
\fancyhf{}
\fancyhead[L]{\small\color{litblue}Froth — Apuntes de ML}
\fancyhead[R]{\small\thepage}
\renewcommand{\headrulewidth}{0.4pt}
\usepackage[colorlinks=true,linkcolor=litblue,urlcolor=litblue]{hyperref}

% Cabecera de cada apunte: \cabecera{numero}{titulo}
\newcommand{\cabecera}[2]{%
  \begin{tcolorbox}[colback=litblue,coltext=white,colframe=litblue,arc=2mm]
    {\Large\bfseries Apunte #1 — #2}\\[3pt]
    {\small Proyecto Froth · aprender Machine Learning construyendo}
  \end{tcolorbox}\vspace{0.8em}}

% Cajas temáticas
\newtcolorbox{concepto}[1]{colback=blue!4,colframe=litblue,title={Concepto: #1},arc=1.5mm}
\newtcolorbox{practica}{colback=green!5,colframe=litgreen,title={Pruébalo tú (teclado en mano)},arc=1.5mm}
\newtcolorbox{ojo}{colback=orange!8,colframe=litorange,title={Ojo / error típico},arc=1.5mm}
\newtcolorbox{resumen}{colback=gray!8,colframe=gray!60,title={En una frase},arc=1.5mm}

\newcommand{\codigo}[1]{\texttt{#1}}

\begin{document}
\cabecera{15}{Ponerle nombre a las burbujas: c-TF-IDF}

\section{El problema}

HDBSCAN te da grupos llamados ``0'', ``1'', ``2''. Para el mapa necesitamos nombres con
significado. ¿Cómo se nombra un grupo de 49 papers \emph{automáticamente}?

\section{Por qué NO sirven las palabras más frecuentes}

Si tomas las palabras más repetidas de cada cluster, en TODOS sale lo mismo: ``mineral'',
``flotation'', ``study'', ``results''. Frecuente $\neq$ característico. Lo que distingue a un
grupo son las palabras que él usa mucho \textbf{y los demás poco}.

\begin{concepto}{TF-IDF y su versión por clases (c-TF-IDF)}
\textbf{TF-IDF} puntúa cada palabra combinando dos factores: qué tan frecuente es en un
documento (TF) y qué tan \emph{rara} es en el resto (IDF). Palabra frecuente aquí + rara en
los demás = puntuación alta = característica.

La variante \textbf{c-TF-IDF} (class-based) aplica esto a clusters: se pegan todos los
abstracts de cada cluster como si fueran \textbf{un solo mega-documento}, y se comparan los
mega-documentos entre sí. Las palabras top de cada mega-documento son la etiqueta del cluster.
Es el mismo truco que usa BERTopic por dentro.
\end{concepto}

\section{Lo que salió con TUS clusters}

\begin{center}
\begin{tabular}{lp{5.2cm}p{5.2cm}}
\textbf{Grupo} & \textbf{Keywords (top c-TF-IDF)} & \textbf{Lectura} \\
\hline
Cluster 0 (49) & melt, crystallization, Nb, Zr, Hf & Geología/petrología de pegmatitas y
granitos de metales raros \\
Cluster 2 (43) & flotation, collector, froth, lepidolite & Flotación de lepidolita:
colectores y espumas \\
Cluster 1 (25) & commodity, old scrap, USGS, women, men & Economía/reciclaje del litio + los
papers off-topic (cuarentena) \\
\end{tabular}
\end{center}

El mapa tiene sentido mineralógico completo: a un lado la geología del yacimiento, al otro el
procesamiento (flotación), y aparte ``el resto del mundo''.

\begin{ojo}
\textbf{Hallazgo de calidad de datos:} entre las keywords del cluster de flotación se colaron
palabras en \emph{francés} (``du'', ``récupération''). Causa: el corpus tiene abstracts en
francés (¡Beauvoir está en Francia!) y nuestra lista de \emph{stopwords} (palabras vacías que
se ignoran: the, of, and...) solo cubre inglés. No rompe nada, pero ensucia etiquetas.
Arreglos posibles: stopwords multilingües o detectar idioma en la limpieza de Fase 1.
Lección: los datos reales siempre traen una sorpresa más.
\end{ojo}

\begin{concepto}{el algoritmo propone, el experto valida}
Las etiquetas automáticas son un \emph{borrador}: ``melt--crystallization--Nb'' es útil, pero
tú como experto la refinarás a ``Petrogénesis de granitos de metales raros''. Esa división de
trabajo (máquina propone, humano cura) es exactamente cómo se usan estas herramientas en
investigación seria.
\end{concepto}

\begin{resumen}
c-TF-IDF nombra cada cluster con sus palabras características (frecuentes dentro, raras
fuera), pegando cada cluster en un mega-documento. Tus 3 subtemas: geología de pegmatitas,
flotación de lepidolita, y economía/reciclaje (con los off-topic en cuarentena).
\end{resumen}

\end{document}
